
\let\negmedspace\undefined
\let\negthickspace\undefined
\documentclass[journal,12pt,onecolumn]{IEEEtran}
\usepackage{cite}
\usepackage{amsmath,amssymb,amsfonts,amsthm}
\usepackage{algorithmic}
\usepackage{graphicx}
\graphicspath{{./figs/}}
\usepackage{textcomp}
\usepackage{xcolor}
\usepackage{txfonts}
\usepackage{listings}
\usepackage{enumitem}
\usepackage{mathtools}
\usepackage{gensymb}
\usepackage{comment}
\usepackage{caption}
\usepackage[breaklinks=true]{hyperref}
\usepackage{tkz-euclide} 
\usepackage{listings}
\usepackage{gvv}                                        
%\def\inputGnumericTable{}                                 
\usepackage[latin1]{inputenc}     
\usepackage{xparse}
\usepackage{color}                                            
\usepackage{array}
\usepackage{longtable}                                       
\usepackage{calc}                                             
\usepackage{multirow}
\usepackage{multicol}
\usepackage{hhline}                                           
\usepackage{ifthen}                                           
\usepackage{lscape}
\usepackage{tabularx}
\usepackage{array}
\usepackage{float}
\newtheorem{theorem}{Theorem}[section]
\newtheorem{problem}{Problem}
\newtheorem{proposition}{Proposition}[section]
\newtheorem{lemma}{Lemma}[section]
\newtheorem{corollary}[theorem]{Corollary}
\newtheorem{example}{Example}[section]
\newtheorem{definition}[problem]{Definition}
\newcommand{\BEQA}{\begin{eqnarray}}
\newcommand{\EEQA}{\end{eqnarray}}
\newcommand{\define}{\stackrel{\triangle}{=}}
\theoremstyle{remark}
\newtheorem{rem}{Remark}

\begin{document}

\title{5.5.9}
\author{EE25BTECH11057 - Tammineedi Rushil Shanmukha Srinivas}
\maketitle
\renewcommand{\thefigure}{\theenumi}
\renewcommand{\thetable}{\theenumi}
\textbf{Question:} Using elementary row transformations find the inverse of the matrix 
\begin{align}
\myvec{3&0&-1\\
       2&3&0\\
       0&4&1}
\end{align}
\textbf{Solution:}
Given, \begin{align} 
\vec{A}=\myvec{3&0&-1\\
       2&3&0\\
       0&4&1}
       \end{align}
Forming the Augmented Matrix
\begin{align}
\myvec{3&0&-1 &\vrule& 1&0&0 \\
       2&3&0 &\vrule&  0&1&0 \\
       0&4&1 &\vrule&  0&0&1 }
\end{align}
Applying elementary row operations to find the inverse
\begin{align}
\myvec{3&0&-1 &\vrule& 1&0&0 \\
       2&3&0 &\vrule&  0&1&0 \\
       0&4&1 &\vrule&  0&0&1 } \xrightarrow{R_3 \longrightarrow 3R_3-4R_2}
       \myvec{3&0&-1 &\vrule& 1&0&0 \\
       2&3&0 &\vrule&  0&1&0 \\
       -8&0&3 &\vrule&  0&-4&3} \xrightarrow{R_3 \longrightarrow 3R_3+8R_1}
       \myvec{3&0&-1 &\vrule& 1&0&0 \\
       2&3&0 &\vrule&  0&1&0 \\
       0&0&1 &\vrule&  8&-12&9} 
\end{align}
\begin{align}
 \myvec{3&0&-1 &\vrule& 1&0&0 \\
       2&3&0 &\vrule&  0&1&0 \\
       0&0&1 &\vrule&  8&-12&9} \xrightarrow{R_2 \longrightarrow 3R_2-2R_1}
    \myvec{3&0&-1 &\vrule& 1&0&0 \\
       0&9&2 &\vrule&  -2&3&0 \\
       0&0&1 &\vrule&  8&-12&9} \xrightarrow{R_2 \longrightarrow R_2-2R_3}
        \myvec{3&0&-1 &\vrule& 1&0&0 \\
       0&9&0 &\vrule&  -18&27&-18 \\
       0&0&1 &\vrule&  8&-12&9} 
\end{align}
\begin{align}
 \myvec{3&0&-1 &\vrule& 1&0&0 \\
       0&9&0 &\vrule&  -18&27&-18 \\
       0&0&1 &\vrule&  8&-12&9} \xrightarrow{R_2 \longrightarrow \frac{R_2}{9}}
        \myvec{3&0&-1 &\vrule& 1&0&0 \\
       0&1&0 &\vrule&  -2&3&-2 \\
       0&0&1 &\vrule&  8&-12&9} \xrightarrow{R_1 \longrightarrow R_1+R_3}
        \myvec{3&0&0 &\vrule& 9&-12&9 \\
       0&1&0 &\vrule&  -2&3&-2 \\
       0&0&1 &\vrule&  8&-12&9} 
\end{align}
\begin{align}
\myvec{3&0&0 &\vrule& 9&-12&9 \\
       0&1&0 &\vrule&  -2&3&-2 \\
       0&0&1 &\vrule&  8&-12&9} \xrightarrow{R_1 \longrightarrow \frac{R_1}{3}}
       \myvec{1&0&0 &\vrule& 3&-4&3 \\
       0&1&0 &\vrule&  -2&3&-2 \\
       0&0&1 &\vrule&  8&-12&9} 
\end{align}
The right side part of the Augmented Matrix is $\vec{A}^{-1}$
\begin{align}
\vec{A}^{-1} = \myvec{3&-4&3\\
                      -2&3&-2\\
                      8&-12&9}
\end{align}
\end{document}